% NESTED VERSION -----------------------------------------------------------------------------------
\begin{work}
{Software Team Member}
{Robotics @ Maryland}
{College Park, MD}
{Jan 2024 - May 2024}

\item \workdetail{Robotic Arm Workspace Analysis} {Created a tool to assist with optimizing the
      kinematics of a robotic arm by developing a program to plot the 3D workspace, resulting in a
      5\% increase in dexterous workspace volume for the required tasks.}
      %
      {https://github.com/TerraformersURC/workspace-plotter}

      \begin{workdetail-list}

      \item Written C++ code utilizing KDL and Matplot++ for simulating forward kinematics using the
            Monte Carlo method.

      \item Implemented a parser to read the DH parameters from a YAML file using yaml-cpp to
            rapidly test different setups.

      \end{workdetail-list}

% \item Facilitated simulation of a rover design in Fusion 360 using Gazebo and RViz by creating a
%       URDF description package.

% \item Maintained the project repositories and documentation on GitHub.

\end{work}

% % NORMAL VERSION -----------------------------------------------------------------------------------
% \begin{work}
%   {Software Team Member}
%   {Robotics @ Maryland - Student Club}
%   {}
%   {Jan 2024 - May 2024}

%   \item Created a tool to assist with optimizing the kinematics of a robotic arm by developing a
%         C++ program with Orocos-KDL and Matplot++ libraries, to plot the 3D workspace, resulting in
%         a 5\% increase in dexterous volume for the required tasks.

%   % \item Built a URDF description package of a rover from the Fusion 360 model for simulations using
%   %       the ROS ecosystem.

% \end{work}

% % SPICED VERSION ---------------------------------------------------------------------------------
% \begin{work}
%   {Software Team Member}
%   {Robotics @ Maryland}
%   {College Park, MD}
%   {Jan 2024 - May 2024}

%   \item Developed a workspace evaluation tool for serial manipulators, enabling kinematic
%         optimization through 3D simulation and visualization of reachable zones.

%   \item Applied first-principles reasoning to optimize arm joint configurations and enhance
%         dexterous reach for field robotics applications.

% \end{work}
